Sobre cada \texttt{ExamPage} reconocida, los correctores aplican
anotaciones de tres tipos: \textbf{texto enriquecido}, \textbf{trazos a
lápiz} (capturados como \acs{json} con la geometría del trazo y
guardados \textit{inline} en la base de datos) y \textbf{audio}
(almacenado en MinIO y referenciado desde la anotación). Cada anotación
está vinculada a una zona geométrica concreta de la página y a un
problema, lo que permite reconstruir la corrección en su contexto
visual.

La calificación de cada problema admite dos modos complementarios:
\textbf{calificación por rúbrica}, en la que el corrector marca los
criterios cumplidos y el sistema deriva la puntuación del problema a
partir de los pesos definidos, y \textbf{calificación manual}, en la
que el corrector introduce la puntuación numérica directamente. Ambos
modos pueden coexistir sobre el mismo examen, configurándose problema
a problema. La nota global de la \texttt{ExamInstance} se calcula
automáticamente como agregación de las puntuaciones por problema.
